\clearpage
\VerdierPhysicalPageAnchor{0038}%
\VerdierSemanticAnchor{ch1-1-2-4-isomorphisms-cont}
と同型

\[
  m(g)\colon F_eT(g)\longrightarrow T'(g)F_e,
  \qquad g\in G,
\]
との組であって、図式 (1.2.3.2) が可換となるものである。

\VerdierSemanticAnchor{ch1-1-2-4-characterization}
したがって $G^\circ$-関手とは「$G^\circ$ の作用と同型を除いて可換する」関手
である。これを $(F_e,m)$ と書く。

\VerdierSemanticAnchor{ch1-1-2-4-strict}
各 $g\in G$ に対して同型 $m(g)$ が恒等同型であるとき（特に各 $g\in G$ に
対して $F_eT(g)=T'(g)F_e$ となる）、この $G^\circ$-関手を\emph{狭義}という。

\VerdierSemanticAnchor{ch1-1-2-5}
\noindent\textbf{1.2.5.}\quad
$\mathcal C$, $\mathcal C'$, $\mathcal C''$ を三つの次数付き圏とし、
\[
  \mathcal C\xrightarrow{F}\mathcal C'
  \xrightarrow{F'}\mathcal C'',\qquad F''=F'F,
\]
を三つの次数付き関手とする。$\mathcal C$, $\mathcal C'$, $\mathcal C''$ に
正規化された切断系を選ぶ。以上から三つの $G^\circ$-圏
$(\mathcal C_e,G^\circ,T,c)$,
$(\mathcal C'_e,G^\circ,T',c')$,
$(\mathcal C''_e,G^\circ,T'',c'')$ と、三つの $G^\circ$-関手
$(F_e,m)$, $(F'_e,m')$, $(F''_e,m'')$ を得る。直ちに次の関係が分かる：

\VerdierSemanticAnchor{ch1-1-2-5-eq1}
\begin{equation}
  \begin{cases}
    F''_e=F'_eF_e,\\[0.5em]
    m''(g)=\bigl(m'(g)\star F_e\bigr)
      \bigl(F'_e\star m(g)\bigr),\qquad g\in G.
  \end{cases}
  \tag{1.2.5.1}
\end{equation}

\VerdierSemanticAnchor{ch1-1-2-6}
\noindent\textbf{定義 1.2.6.}\quad
$(\mathcal C_e,G^\circ,T,c)$,
$(\mathcal C'_e,G^\circ,T',c')$,
$(\mathcal C''_e,G^\circ,T'',c'')$ を三つの $G^\circ$-圏とし、

\[
\begin{tikzcd}[column sep=large]
(\mathcal C_e,G^\circ,T,c)
  \arrow[r,"{(F_e,m)}"]
&
(\mathcal C'_e,G^\circ,T',c')
  \arrow[r,"{(F_e^{\prime},m^{\prime})}"]
&
(\mathcal C''_e,G^\circ,T'',c'')
\end{tikzcd}
\]
を二つの $G^\circ$-関手とする。

\VerdierSemanticAnchor{ch1-1-2-6-composite}
等式 (1.2.5.1) は、$G^\circ$-関手 $(F_e,m)$ と $(F'_e,m')$ の
\emph{合成関手}と呼ばれる $G^\circ$-関手を定義する。

\VerdierSemanticAnchor{ch1-1-2-7}
\noindent\textbf{1.2.7.}\quad
$F,F'$ を二つの $G$-次数付き圏 $\mathcal C,\mathcal C'$ の間の次数付き関手、
$u\colon F\to F'$ を次数 $e$ の関手の射とする。$\mathcal C,\mathcal C'$ に
正規化された切断系を選ぶ。二つの $G^\circ$-圏
$(\mathcal C_e,G^\circ,T,c)$, $(\mathcal C'_e,G^\circ,T',c')$ と、二つの
$G^\circ$-関手 $(F_e,m)$, $(F'_e,m')$ を得る。射 $u$ は関手の射

\VerdierSemanticAnchor{ch1-1-2-7-eq1}
\begin{equation}
  u_e\colon F_e\longrightarrow F'_e
  \tag{1.2.7.1}
\end{equation}
を定め、次の図式が可換となる：

\VerdierSemanticAnchor{ch1-1-2-7-eq2}
\begin{equation}
\begin{tikzcd}[column sep=huge,row sep=large]
F_eT(g)
  \arrow[r,"u_e\star T(g)"]
  \arrow[d,"m(g)",swap]
&
F_e^{\prime}T(g)
  \arrow[d,"m^{\prime}(g)"]
\\
T^{\prime}(g)F_e
  \arrow[r,"T^{\prime}(g)\star u_e",swap]
&
T^{\prime}(g)F_e^{\prime}\mathpunct{.}
\end{tikzcd}
\tag{1.2.7.2}
\end{equation}
